\section{正交增强与全局设置}
\label{sec:appendix}

\paragraph{正交增强。}
multi-crop 聚合与 pixel-to-image 融合是标准的空间/分数聚合手段，与本文“逐图正常参照”机制正交。它们能进一步提升数值，但不改变本文的机制结论；为保持正文聚焦，正文不将其计入核心方法。

\paragraph{全局设置 vs 逐数据集调参。}
\tabref{tab:global} 区分统一全局设置与逐数据集调参上界。我们主张以全局设置作为公平结论，调参行仅作为上界参考，以避免“测试集调参”的质疑。

\begin{table}[h]
\centering
\caption{全局设置 vs 逐数据集调参上界（pixel AUPRO，\%）。目标占位数值，非当前结果口径。}
\label{tab:global}
\small
\begin{tabular}{lcc}
\toprule
数据集 & 全局设置 & 逐数据集调参（上界） \\
\midrule
MPDD & 88.4 & 90.5 \\
BTAD & 79.5 & 78.5 \\
DTD-Synth & 90.7 & 91.5 \\
\bottomrule
\end{tabular}
\end{table}

\paragraph{可复现性说明。}
本文为目标占位数值，用于锁定早期叙事与判定标准；当前版本应以已复现结果表替换，并从源论文核实所有外部对比数值与引用。
